%=====================================================================
% PROJETO RESIDENCIAL I — CASA TERREA COM DUAS SUITES
% Base eletrica: ABNT NBR 5410 + NT.00001.EQTL Rev.09
% Planta-base didatica com logica arquitetonica e mobiliario.
% Atualizacao: agosto de 2026
%=====================================================================

\ifdefined\ArqParedeExt\else\input{projetos/base/plantas_arquitetonicas.tex}\fi

\section{Projeto Residencial I — casa térrea com duas suítes}

% Símbolos elétricos locais
\providecommand{\PRLuz}[2]{\begin{scope}[shift={(#1,#2)}]\draw[corAccent,thick] (0,0) circle (0.14);\draw[corAccent,thick] (-0.10,-0.10)--(0.10,0.10);\draw[corAccent,thick] (-0.10,0.10)--(0.10,-0.10);\end{scope}}
\providecommand{\PRTUG}[2]{\begin{scope}[shift={(#1,#2)}]\draw[corPrim,thick] (0,0) circle (0.13);\draw[corPrim,thick] (-0.16,0)--(0.16,0);\end{scope}}
\providecommand{\PRTUE}[3]{\begin{scope}[shift={(#1,#2)}]\draw[corAt,thick] (0,0) circle (0.17);\node[font=\tiny\bfseries,corAt] at (0,0){#3};\end{scope}}
\providecommand{\PRQD}[3]{\begin{scope}[shift={(#1,#2)}]\draw[corPrim,fill=corDef!10,thick,rounded corners=1pt] (-0.38,-0.22) rectangle (0.38,0.22);\node[font=\tiny\bfseries] at (0,0){#3};\end{scope}}

%---------------------------------------------------------------------
% Planta 15,00 x 11,00 m. Frente em y=0.
% Ala íntima à esquerda; social ao centro; garagem/serviço à direita.
%---------------------------------------------------------------------
\newcommand{\CasaResidencialBase}{%
  % contorno
  \draw[arqExt] (0,0) rectangle (15,11);
  % ala íntima e hall
  \ArqParedeInt{5.2}{0}{5.2}{11}
  \ArqParedeInt{6.4}{0}{6.4}{11}
  \ArqParedeInt{0}{3.6}{5.2}{3.6}
  \ArqParedeInt{0}{7.2}{5.2}{7.2}
  % banheiros internos / social
  \ArqParedeInt{3.5}{0}{3.5}{2.0}
  \ArqParedeInt{3.5}{2.0}{5.2}{2.0}
  \ArqParedeInt{3.5}{4.8}{5.2}{4.8}
  \ArqParedeInt{3.5}{4.8}{3.5}{7.2}
  \ArqParedeInt{3.5}{8.4}{5.2}{8.4}
  \ArqParedeInt{3.5}{8.4}{3.5}{11}
  % social/cozinha
  \ArqParedeInt{6.4}{6.4}{11.2}{6.4}
  % garagem / serviço / varanda gourmet
  \ArqParedeInt{11.2}{0}{11.2}{11}
  \ArqParedeInt{11.2}{5.6}{15}{5.6}
  \ArqParedeInt{11.2}{8.2}{15}{8.2}

  % acessos: principal diretamente na sala; garagem para sala; serviço externo
  \ArqPortaHUp{7.25}{0}{1.10}
  \draw[white,line width=5pt] (11.2,2.2)--(11.2,3.2);\draw[black!60,thin] (11.2,2.2)--(10.2,2.2);
  \ArqPortaVLeft{15}{6.25}{0.90}
  \draw[white,line width=5pt] (11.2,9.1)--(11.2,10.7);\draw[arqGlass] (11.2,9.1)--(11.2,10.7);
  % hall interno aberto para a sala
  \draw[white,line width=5pt] (6.4,2.45)--(6.4,3.65);
  % portas dos quartos/banheiros para hall
  \ArqPortaVLeft{5.2}{2.45}{0.90}
  \ArqPortaVLeft{5.2}{3.85}{0.90}
  \ArqPortaVLeft{5.2}{7.45}{0.90}
  \ArqPortaVLeft{5.2}{0.55}{0.80}
  % banheiros das suítes: acesso exclusivo por dentro do quarto
  \ArqPortaVLeft{3.5}{5.05}{0.80}
  \ArqPortaVLeft{3.5}{8.65}{0.80}
  % cozinha -> serviço
  \ArqPortaVRight{11.2}{6.75}{0.85}
  % vão de garagem
  \draw[white,line width=5.5pt] (11.55,0)--(14.70,0);\draw[black!30,very thin] (11.55,0.10)--(14.70,0.10);

  % janelas
  \ArqJanV{0.45}{2.75}{0}
  \ArqJanV{4.15}{6.55}{0}
  \ArqJanV{7.75}{10.20}{0}
  \ArqJanH{8.0}{10.2}{0}
  \ArqJanH{7.3}{9.8}{11}
  \ArqJanH{0.7}{2.6}{11}
  \ArqJanH{12.0}{14.1}{11}

  % mobiliário - quarto 3
  \ArqCama{0.55}{0.45}{2.15}{2.65}
  \ArqArmario{2.85}{2.45}{0.55}{0.95}
  \node[arqLabel] at (1.75,3.25){QUARTO 3};
  % banho social
  \ArqChuveiro{3.65}{0.15}{1.25}{0.85}\ArqVaso{4.55}{1.35}\ArqPia{3.62}{1.20}{0.65}{0.45}
  \node[arqSub] at (4.30,1.82){BANHO SOCIAL};

  % suíte 2
  \ArqCama{0.55}{4.05}{2.20}{2.65}\ArqArmario{2.85}{3.95}{0.50}{1.05}
  \ArqChuveiro{3.65}{6.05}{1.25}{0.95}\ArqVaso{4.55}{5.35}\ArqPia{3.62}{5.05}{0.65}{0.45}
  \node[arqLabel] at (1.75,6.85){SUÍTE 2};\node[arqSub] at (4.35,6.90){B2};

  % suíte master
  \ArqCama{0.55}{7.75}{2.35}{2.70}\ArqArmario{2.95}{7.45}{0.45}{0.85}
  \ArqChuveiro{3.65}{10.00}{1.25}{0.85}\ArqVaso{4.55}{9.20}\ArqPia{3.62}{8.65}{0.65}{0.45}
  \node[arqLabel] at (1.85,10.65){SUÍTE MASTER};\node[arqSub] at (4.35,10.70){B1};
  \node[arqSub,rotate=90] at (5.80,5.50){HALL ÍNTIMO};

  % sala/jantar
  \ArqSofa{7.00}{0.85}{3.10}{1.05}\ArqBancada{10.50}{0.95}{0.35}{2.25}
  \ArqMesa{7.25}{4.20}{2.55}{1.10}
  \node[arqLabel] at (8.80,3.20){SALA / JANTAR};
  % cozinha
  \ArqBancada{6.75}{10.15}{3.50}{0.55}\ArqBancada{10.35}{7.05}{0.55}{3.65}
  \ArqFogao{9.25}{10.05}\ArqGeladeira{10.15}{6.65}\ArqBancada{7.40}{7.55}{2.30}{0.80}
  \node[arqLabel] at (8.60,9.45){COZINHA};
  % garagem
  \ArqCarro{12.00}{0.65}{2.15}{4.25}\node[arqLabel] at (13.10,5.20){GARAGEM};
  % serviço
  \ArqLavadora{12.00}{6.05}\ArqBancada{13.0}{6.0}{1.45}{0.55}\node[arqLabel] at (13.10,7.65){SERVIÇO};
  % gourmet / quintal coberto
  \ArqBancada{12.0}{9.85}{2.35}{0.55}\ArqMesa{12.20}{8.70}{1.55}{0.75}\node[arqLabel] at (13.10,10.75){VARANDA GOURMET};

  \CotaH{0}{15}{11.18}{15,00 m}\CotaV{0}{11}{15.18}{11,00 m}
}

\begin{frame}{Residencial I — planta arquitetônica funcional}
\centering
\begin{tikzpicture}[scale=0.49,every node/.style={font=\tiny}]
\CasaResidencialBase
\end{tikzpicture}
\vspace{0.5mm}
\scriptsize A entrada principal chega à sala. O hall é interno. A casa possui três banheiros completos: banho social, B2 e B1; estes dois últimos são acessados exclusivamente pelas suítes.
\end{frame}

\begin{frame}{Residencial I — iluminação e circuitos C1/C2}
\centering
\begin{tikzpicture}[scale=0.49,every node/.style={font=\tiny}]
\CasaResidencialBase
\PRQD{6.10}{3.10}{QD1}
\foreach \x/\y in {1.7/2.1,4.25/1.0,1.75/5.4,4.25/5.9,1.8/9.2,4.25/9.6,5.8/5.5,8.8/2.8,8.7/5.2,8.6/8.8,13.1/3.0,13.1/6.9,13.1/9.6}{\PRLuz{\x}{\y}}
% C1: sala, cozinha, servico e gourmet
\draw[corAccent!80!black,dashed,thick] (6.10,3.10)--(8.8,2.8)--(8.7,5.2)--(8.6,8.8)--(13.1,9.6);
\draw[corAccent!80!black,dashed,thick] (8.6,8.8)--(13.1,6.9);
% C2: ala intima, banheiros, hall e garagem
\draw[corAccent!55!black,dashed,thick] (6.10,3.10)--(5.8,5.5)--(1.75,5.4)--(4.25,5.9)--(1.8,9.2)--(4.25,9.6);
\draw[corAccent!55!black,dashed,thick] (6.10,3.10)--(1.7,2.1)--(4.25,1.0);
\draw[corAccent!55!black,dashed,thick] (6.10,3.10)--(13.1,3.0);
\node[font=\tiny\bfseries,fill=white,inner sep=1pt,corAccent!80!black] at (9.5,7.4){C1};
\node[font=\tiny\bfseries,fill=white,inner sep=1pt,corAccent!55!black] at (3.0,5.4){C2};
\end{tikzpicture}
\vspace{0.4mm}
\scriptsize Previsão calculada: C1 = 1,06 kVA e C2 = 1,40 kVA. Os 13 símbolos são pontos de luz; os 2,46 kVA resultam das áreas e dos acréscimos completos de 4 m$^2$, não da quantidade de símbolos.
\end{frame}

\begin{frame}{Residencial I — 43 TUG e oito cargas específicas}
\centering
\begin{tikzpicture}[scale=0.49,every node/.style={font=\tiny}]
\CasaResidencialBase
\PRQD{6.10}{3.10}{QD1}
% C3: sala e dormitorios — 17 pontos
\foreach \x/\y in {0.25/0.35,0.25/3.30,3.25/0.35,3.25/3.30,0.25/3.85,0.25/6.95,3.25/3.85,3.25/6.95,0.25/7.45,0.25/10.75,3.25/7.45,3.25/10.75,6.65/0.35,10.95/0.35,10.95/3.60,10.95/6.15,6.65/6.15}{\PRTUG{\x}{\y}}
% C4: hall e garagem — 9 pontos
\foreach \x/\y in {6.25/0.70,6.25/2.10,6.25/4.20,6.25/7.80,6.25/10.20,11.35/0.80,11.35/4.80,14.65/1.50,14.65/4.80}{\PRTUG{\x}{\y}}
% C5 cozinha; C6 servico; C7 gourmet; C8 banheiros
\foreach \x/\y in {6.70/10.75,8.00/10.75,9.30/10.75,10.95/7.00,10.95/8.50,10.95/10.00}{\PRTUG{\x}{\y}}
\foreach \x/\y in {11.35/5.85,14.65/5.85,11.35/7.95,14.65/7.95}{\PRTUG{\x}{\y}}
\foreach \x/\y in {11.35/8.45,14.65/8.45,11.35/10.75,14.65/10.75}{\PRTUG{\x}{\y}}
\foreach \x/\y in {5.00/1.80,4.95/5.00,4.95/8.65}{\PRTUG{\x}{\y}}
% TUE C9--C16
\PRTUE{4.35}{0.55}{C9}\PRTUE{4.35}{6.55}{C10}\PRTUE{4.35}{10.45}{C11}
\PRTUE{9.60}{10.45}{C12}\PRTUE{12.35}{6.40}{C13}
\PRTUE{0.45}{2.55}{C14}\PRTUE{0.45}{6.25}{C15}\PRTUE{0.45}{9.95}{C16}
\end{tikzpicture}
\vspace{0.3mm}
\tiny C3 sala/quartos (17) · C4 hall/garagem (9) · C5 cozinha (6) · C6 serviço (4) · C7 gourmet (4) · C8 banheiros (3) · C9--C11 chuveiros · C12 forno · C13 lavadora · C14--C16 ar-condicionado.
\end{frame}

\begin{frame}[shrink=7]{Residencial I — memória dos mínimos por ambiente}
\scriptsize
\begin{tabularx}{\textwidth}{>{\bfseries}p{2.55cm}c c c c >{\centering\arraybackslash}X}
\toprule
Ambiente & Área (m$^2$) & Perím. (m) & Luz (VA) & TUG & TUG (VA)\\
\midrule
Quarto 3 & 15,3 & 17,6 & 220 & 4 & 400\\
Suíte 2 & 14,6 & 17,6 & 220 & 4 & 400\\
Suíte master & 15,3 & 18,0 & 220 & 4 & 400\\
Banheiros (3) & 3,4/4,1/4,4 & critério próprio & 300 & 3 & 1.800\\
Hall íntimo & 13,2 & 24,4 & 160 & 5 & 500\\
Sala/jantar & 30,7 & 22,4 & 460 & 5 & 500\\
Cozinha & 22,1 & 18,8 & 340 & 6 & 2.100\\
Garagem & 21,3 & 18,8 & 280 & 4 & 400\\
Serviço & 9,9 & 12,8 & 100 & 4 & 1.900\\
Varanda gourmet$^*$ & 10,6 & 13,2 & 160 & 4 & 1.900\\
\midrule
\textbf{Total} & -- & -- & \textbf{2.460} & \textbf{43} & \textbf{10.300}\\
\bottomrule
\end{tabularx}
\begin{caixaObs}
Áreas e perímetros seguem a geometria interna desenhada; os dormitórios têm forma em L. $^*$A varanda gourmet foi tratada como ambiente análogo a cozinha por possuir bancada de preparo.
\end{caixaObs}
\end{frame}

\begin{frame}[shrink=6]{Residencial I — circuitos C1 a C8}
\scriptsize
\begin{tabularx}{\textwidth}{c p{3.25cm} c c c c c X}
\toprule
C & Uso & $S$ & $I_b$ & Cu & DJ & Fase & Solução adotada\\
\midrule
C1 & Luz social/cozinha/serviço/gourmet & 1,06 kVA & 4,8 A & 1,5 mm$^2$ & 10 A & A & RCBO 30 mA\\
C2 & Luz íntima/banhos/hall/garagem & 1,40 kVA & 6,4 A & 1,5 mm$^2$ & 10 A & B & RCBO 30 mA\\
C3 & TUG sala e dormitórios & 1,70 kVA & 7,7 A & 2,5 mm$^2$ & 16 A & A & RCBO 30 mA\\
C4 & TUG hall e garagem & 0,90 kVA & 4,1 A & 2,5 mm$^2$ & 16 A & B & RCBO 30 mA\\
C5 & TUG cozinha & 2,10 kVA & 9,5 A & 2,5 mm$^2$ & 16 A & C & RCBO 30 mA\\
C6 & TUG serviço & 1,90 kVA & 8,6 A & 2,5 mm$^2$ & 16 A & C & RCBO 30 mA\\
C7 & TUG gourmet & 1,90 kVA & 8,6 A & 2,5 mm$^2$ & 16 A & A & RCBO 30 mA\\
C8 & TUG banheiros & 1,80 kVA & 8,2 A & 2,5 mm$^2$ & 16 A & C & RCBO 30 mA\\
\bottomrule
\end{tabularx}
\begin{caixaObs}
RCBO individuais evitam neutros compartilhados a jusante e limitam a área afetada. Tipo A é a premissa inicial; confirmar forma de onda residual, corrente de fuga e fabricante.
\end{caixaObs}
\end{frame}

\begin{frame}[shrink=6]{Residencial I — circuitos C9 a C16}
\scriptsize
\begin{tabularx}{\textwidth}{c p{3.15cm} c c c c c X}
\toprule
C & Equipamento & $P$ & $I_b$ & Cu & DJ & Fase & Proteção/observação\\
\midrule
C9 & Chuveiro banho social & 4,50 kW & 20,5 A & 4 mm$^2$ & 25 A & A & dedicado; RCBO 30 mA\\
C10 & Chuveiro suíte 2 & 4,50 kW & 20,5 A & 4 mm$^2$ & 25 A & B & dedicado; RCBO 30 mA\\
C11 & Chuveiro suíte master & 4,50 kW & 20,5 A & 4 mm$^2$ & 25 A & C & dedicado; RCBO 30 mA\\
C12 & Forno elétrico & 3,50 kW & 15,9 A & 4 mm$^2$ & 20 A & B & dedicado; RCBO 30 mA\\
C13 & Lavadora & 1,20 kW & 5,5 A & 2,5 mm$^2$ & 16 A & A & dedicado; RCBO 30 mA\\
C14--C16 & Ar-condicionado quartos & 3$\times$1,00 kW & 4,5 A cada & 2,5 mm$^2$ & 16 A & C/B/A & dedicados; fabricante\\
\bottomrule
\end{tabularx}
\begin{caixaAt}
As seções e proteções são seleções iniciais. Antes do executivo: método de instalação, agrupamento, temperatura, $I_z$, queda, curto máximo/mínimo, capacidade de interrupção e manual de cada equipamento.
\end{caixaAt}
\end{frame}

\begin{frame}[shrink=4]{Residencial I — entrada e balanceamento fechados}
\small
\begin{columns}[T,onlytextwidth]
\column{0.49\textwidth}
\begin{caixaProjeto}[title=Balanceamento por valores instalados]
\begin{align*}
S_A&=\SI{11.36}{kVA}\quad (\SI{51.6}{A})\\
S_B&=\SI{11.30}{kVA}\quad (\SI{51.4}{A})\\
S_C&=\SI{11.30}{kVA}\quad (\SI{51.4}{A})
\end{align*}
\scriptsize A: C1+C3+C7+C9+C13+C16\\
B: C2+C4+C10+C12+C15\\
C: C5+C6+C8+C11+C14
\end{caixaProjeto}
\column{0.49\textwidth}
\begin{caixaNorma}[title=Equatorial Maranhão]
Hipótese declarada do estudo: \SI{33.96}{kW} equivalentes. Na Tabela 1 da NT.00001.EQTL Rev.09:
\begin{itemize}
\item faixa 24,1--38 kW;
\item atendimento 380/220 V, 3F+N;
\item disjuntor de entrada 63 A;
\item condutor Cu do cliente no padrão: mínimo 10 mm$^2$.
\end{itemize}
\end{caixaNorma}
\end{columns}
\begin{caixaAt}
Para ilustrar o enquadramento, os VA mínimos de luz/TUG foram tomados numericamente como W. Na ligação real, declarar as potências ativas nominais efetivamente instaladas conforme a definição da NT.00001; demanda não substitui carga instalada nessa tabela. O alimentador interno é recalculado pela NBR 5410.
\end{caixaAt}
\end{frame}

\begin{frame}{Residencial I — unifilar coerente com 16 circuitos}
\centering
\begin{tikzpicture}[scale=0.76,every node/.style={font=\tiny}]
\tikzset{b/.style={draw=corPrim,fill=corDef!7,rounded corners,minimum width=2.05cm,minimum height=0.58cm,align=center}}
\node[b](r)at(0,4.6){Rede EQTL\\380/220 V};
\node[b](m)at(0,3.70){Medição\\padrão de entrada};
\node[b](dg)at(0,2.80){DG 3P 63 A\\+ DPS coordenado};
\node[b](qd)at(0,1.85){QD1\\barras N e PE separadas};
\node[b](g1)at(-4.2,0.45){C1--C8\\luz/TUG + RCBO};
\node[b](g2)at(-1.4,0.45){C9--C11\\chuveiros + RCBO};
\node[b](g3)at(1.4,0.45){C12--C13\\forno/lavadora};
\node[b](g4)at(4.2,0.45){C14--C16\\ar-condicionado};
\node[b,draw=corNorma,fill=corNorma!7](bep)at(5.0,2.65){BEP\\PE + equipot.};
\draw[-{Stealth},thick,corPrim](r)--(m)--(dg)--(qd);
\foreach \x in {g1,g2,g3,g4}{\draw[-{Stealth},thick](qd)--(\x);}
\draw[corNorma,thick](qd)--(bep);
\draw[corNorma,thick](bep.south)--(5.0,1.35)node[ground]{};
\node[font=\tiny,corNorma,anchor=west] at (5.15,1.35){eletrodo(s)};
\end{tikzpicture}
\vspace{1mm}
\scriptsize O PE não é um ``terra isolado'' por circuito; todos os circuitos retornam ao barramento PE, ligado ao BEP e ao sistema comum de aterramento/equipotencialização.
\end{frame}
